\documentclass[10pt,a4paper]{amsart}
\usepackage[margin=19mm]{geometry}
\usepackage{amsmath,amssymb,mathtools}
\usepackage[scr=rsfs,cal=euler]{mathalfa}
\usepackage{xfrac}
\usepackage[unicode,hidelinks,bookmarks=false]{hyperref}
\newcommand{\ie}{i.e.\ }
\newcommand{\cf}{cf.\ }
\newcommand{\resp}{respectively\ }
\newcommand{\Subsection}{\subsection}
\providecommand{\nobreakdash}{\nobreak\hbox{-}}
\setlength{\emergencystretch}{2em}
\multlinegap0pt
% ---------------------------------------------------------------------------
% Editorial AMS wrapper prelude. The theorem environments and the mathematical macros below are
% transcribed from the paper's own preamble (cached-originals/source0.tex lines 22--34); the AMS amsart
% class, the named format packages of the pinned TeX Live runtime and this editorial formatting are not
% part of the author's text. No mathematics is changed.
% ---------------------------------------------------------------------------
\newtheorem{theorem}{Theorem}[section]
\newtheorem{definition}[theorem]{Definition}
\newtheorem{lemma}[theorem]{Lemma}
\newtheorem{corollary}[theorem]{Corollary}
\newtheorem{question}[theorem]{Question}
\newtheorem{example}[theorem]{Example}
\newtheorem{proposition}[theorem]{Proposition}
\newtheorem{remark}[theorem]{Remark}
\newtheorem{assumption}[theorem]{Assumption}
\newtheorem{conjecture}[theorem]{Conjecture}



\begin{document}
\title[The Yang--Baxter Equation and Characteristic Finite Simple Quotients of $F_2$]{The Yang--Baxter Equation and Characteristic Finite Simple Quotients of the Free Group of Rank $2$}
\author{Liam Hanany}
\thanks{The author was supported by the European Research Council (ERC) under the European Union''s Horizon 2020 (N. 882751).}
\date{Source: arXiv:2510.03210v2; distributed under CC BY 4.0}
\maketitle
\noindent\emph{Source, version and licence.} \emph{The Yang-Baxter Equation and Characteristic Finite Simple Quotients of the Free Group of Rank 2}, by Liam Hanany. The source of this editable unit is the e-print of \textbf{version v2} of arXiv:2510.03210, https://arxiv.org/abs/2510.03210v2, whose material is distributed under the Creative Commons Attribution 4.0 International licence, http://creativecommons.org/licenses/by/4.0/, as recorded in the arXiv licence metadata for this paper and shown on that abstract page. The whole of the body below is version v2, the newest version of the paper. Full rights notice: licenses/arxiv-2510.03210-CC-BY-4.0.txt.\par\smallskip
\noindent\textit{Editorial scope.} Counted unit: original Theorem 1.2 of the article, the statement carried by the source environment \texttt{theorem} with source label \texttt{Large rank quotients}, cached-originals/source0.tex lines 65--68, printed as \textbf{Theorem 1.2} exactly as in the article: for every $n \geq 3$ there are infinitely many prime powers $q$ such that there is a characteristic subgroup $C \leq F_2$ with $F_2 / C \simeq \mathrm{PSL}_{\binom{n}{2}}(\mathbb{F}_q)$. It is delivered together with the author\'s own complete proof as the single primary proof body, source0.tex lines 974--988 (15 lines): the \texttt{proof} environment whose optional argument names it "Proof of Theorem \textbackslash ref\{Large rank quotients\}", which runs the Zariski-density and specialization argument for the $B_4$-representation, applies the Strong Approximation Theorem to obtain surjections onto $\mathrm{PSL}_N(\mathbb{F}_p)$ for infinitely many primes, extends the representation to $\mathrm{Aut}(F_2)$ and concludes that the kernel is characteristic. No proof marker is added and no mathematics is written, repaired or re-derived; the author\'s notation and the printed slips of the source are kept literal. Necessary complete context is retained uncounted, in source order: source0.tex lines 257--263, the lemma identifying the semi-direct product structure that produces the $\mathrm{Aut}(F_2)$-action; lines 683--686, the Knizhnik--Zamolodchikov Zariski-density input of the counted proof; lines 799--808, the Zariski-density theorem quoted in the first line of the counted proof; lines 923--928, the corollary extending the representation to $\mathrm{Aut}(F_2)$; and lines 949--959, the two specialization lemmas on good specializations and on the eigenvalues of $\sigma_1\sigma_3^{-1}$ that the counted proof invokes. The article\'s own numbering is reproduced by explicit counter settings: the counter \texttt{theorem}, declared by \texttt{\textbackslash newtheorem\{theorem\}\{Theorem\}[section]} and shared by definition, lemma, corollary, question, example, proposition, remark, assumption and conjecture, is set before each retained passage, so the counted statement prints as Theorem 1.2, the two specialization lemmas print as their original numbered lemmas, and the corollary and the Zariski-density theorem keep their original numbers, exactly as in the article. Every \texttt{\textbackslash ref} and \texttt{\textbackslash eqref} inside every retained passage resolves inside this delivered file. The retained passages contain no \texttt{\textbackslash cite} command at all, so this unit delivers no bibliography entry; the article\'s own bibliography data file \texttt{bibli.bib} is neither spliced into the bundled source nor redistributed. The retained passages contain no figure environment and no graphics-inclusion command, so no artwork is redistributed. The source\'s own class is AMS \texttt{amsart} (\texttt{\textbackslash documentclass[letterpaper]\{amsart\}}); the e-print ships no class or style file, so no publisher class is shipped, used or redistributed, and the wrapper is the same AMS \texttt{amsart} class with the paper\'s own theorem declarations transcribed from its preamble. Every declared body of this unit is an exact, unmodified span of source0.tex: no body is a bounded passage comparison, \texttt{omit} was not used anywhere, and consequently no fidelity receipt is asserted in delivery-evidence.json. The complete concatenated original source is bundled as sources/186345/source0.tex. Omitted from this editable unit and therefore absent from it are source0.tex lines 1--64; 69--256; 264--682; 687--798; 809--922; 929--948; 955--955; 960--973; 989--1172, that is the article\'s own preamble, title block and abstract, the introduction with its citation-dense discussion of the literature, the quandle and $B_4$-action sections, the $\mathrm{PSL}_2(\mathbb{F}_p)$ sections, the proof of the companion alternating-quotients theorem, the large-rank-quotients section, the concluding questions and conjectures, and the article\'s bibliography data, all of which remain present in the bundled source but are not delivered here. The AMS \texttt{amsart} wrapper, the named format packages of the pinned TeX Live runtime and this editorial source, licence and scope paragraph are editorial formatting only.\par\smallskip
\setcounter{section}{1}
\setcounter{equation}{0}
\setcounter{theorem}{1}
\begin{theorem}
\label{Large rank quotients}
    For every $n \geq 3$ there are infinitely many prime powers $q$ such that there is a characteristic subgroup $C \leq F_2$ with $F_2 / C \simeq \mathrm{PSL}_{\binom{n}{2}}(\mathbb{F}_q)$.
\end{theorem}

\setcounter{section}{3}
\setcounter{equation}{1}
\setcounter{theorem}{4}
\begin{lemma}
\label{semi-direct product AutF2}
\begin{enumerate}
\item The map $\varphi: B_4 \to B_4$ defined on generators by $\varphi(\sigma_i) = \sigma_i^{-1}$ is an automorphism. Moreover, restricting the automorphism to $B_4/Z(B_4)$, we get an isomorphism $\mathrm{Aut}(F_2) \simeq B_4/Z(B_4) \rtimes_\varphi \mathbb{Z}/2$.
\item For every $n \geq 2$, the map $\psi: B_n \to B_n$ defined on generators by $\psi(\sigma_i) = \sigma_{n-i}$ is an automorphism. Moreover, this automorphism is inner, and given by conjugation by the Garside element $\Delta_n = (\sigma_1 \cdots \sigma_{n-1}) \cdot (\sigma_1 \cdots \sigma_{n-2}) \cdots (\sigma_1\sigma_2) \cdot \sigma_1.$
\end{enumerate}
\end{lemma}

\setcounter{section}{7}
\setcounter{equation}{0}
\setcounter{theorem}{1}
\begin{theorem}
\label{Knizhinik Zamolodichikov Zariski density}
    The action of $B_n$ on $\textbf{W}_{n,\ell}$ has Zariski-dense image in $\mathrm{PSL}_{\binom{n + \ell - 2}{\ell}}(\bar{\mathbb{K}})$.
\end{theorem}

\setcounter{section}{7}
\setcounter{equation}{0}
\setcounter{theorem}{9}
\begin{theorem}
\label{lie algebra amplification}
    Let $V$ be a finite-dimensional irreducible representation of a connected algebraic group $G$. Let $H \leq G$ be a Zariski-closed, connected subgroup. Let $V|_H = \bigoplus_{k=1}^m W_k$ be the irreducible decomposition of the restriction to $H$. Assume that
    \begin{enumerate}
        \item $(W_k)_{k=1}^m$ are pairwise non-isomorphic and non-dual as projective representations of $H$.
        \item For every $k$, $H$ surjects onto $\mathrm{PSL}(W_k)$.
        \item At most one of the $W_i$ has dimension $1$ and at most one has dimension $2$.
    \end{enumerate}
Then $G$ surjects onto $\mathrm{PSL}(V)$.
\end{theorem}

\setcounter{section}{7}
\setcounter{equation}{0}
\setcounter{theorem}{17}
\begin{corollary}
\label{Extension of representation to AutF2}
There is a matrix $J \in \mathrm{PSL}(\textbf{V}_{4,\ell})$ such that $J\sigma_i^TJ^{-1} = \sigma_i$.
Therefore the automorphism $\varphi$ of $\mathrm{PSL}(\textbf{V}_{4,\ell})$ given by $A \mapsto J(A^T)^{-1}J^{-1}$ satisfies $\varphi \sigma_i \varphi^{-1} = \sigma_i^{-1}$. 
Moreover, $\varphi$ preserves $\mathrm{PSL}(\textbf{W}_{4, \ell})$, and $\varphi|_{\mathrm{PSL}(\textbf{W}_{4, \ell})}$ is an involution.
\end{corollary}

\setcounter{section}{7}
\setcounter{equation}{0}
\setcounter{theorem}{20}
\begin{lemma}
\label{Open subset of good specializations}
Let $A$ be an integral domain, finitely generated over $\mathbb{Q}$, with fraction field $K$. Let $\textbf{G}$ be a connected and absolutely simple linear algebraic group over $K$. Let $\Gamma \leq \mathrm{GL}_n(A)$ and assume that its Zariski-closure over $K$ is $\textbf{G}.$

Denote by $\mathcal{G}$ the Zariski closure of $\Gamma$ over $A$, which has generic fiber $\textbf{G}$. Then there is an open subset $U \subseteq \mathrm{spec}(A)$ such that for every closed point $x \in U$ the image of $\Gamma$ in $\mathcal{G}(k_x)$ is either finite or Zariski-dense, where $k_x$ is the residue field of $\mathrm{spec}(A)$ at $x.$
\end{lemma}

\setcounter{section}{7}
\setcounter{equation}{0}
\setcounter{theorem}{21}
\begin{lemma}
\label{Eigenvalues of sigma1sigma3inv}
Consider the action of $\sigma_1\sigma_3^{-1}$ on $\textbf{W}_{4, \ell}$, for $\ell \geq 1.$ Then all eigenvalues of this matrix have the form $\pm q^as^b$ for some $(a, b) \in \mathbb{Z}^2$, and furthermore these eigenvalues are not all equal up to sign.
\end{lemma}

\setcounter{section}{7}
\setcounter{equation}{0}
\setcounter{theorem}{22}
\begin{proof}[Proof of Theorem \ref{Large rank quotients}]
Let $\ell \geq 1$ and $N = \binom{\ell + 2}{2}$. By Theorem \ref{Knizhinik Zamolodichikov Zariski density}, $B_4 \to \mathrm{PSL}_N(\bar{\mathbb{K}})$ has Zariski-dense image. By Lemma \ref{Eigenvalues of sigma1sigma3inv}, the image of $F_2$ is non-trivial, and therefore as the subgroup $F_2 \lhd B_4$ is normal, it has Zariski-dense image. Let $\pm q^as^b \neq \pm 1$ be a quotient of two eigenvalues of $\sigma_1\sigma_3^{-1}$ that are not equal up to sign.

Since $F_2$ is also contained in the derived subgroup of $B_4$, then the map $B_4 \to \mathrm{PGL}_N(\mathbb{K})$ restricts to a map $F_2 \to \mathrm{PSL}_N(\mathbb{K})$, with Zariski dense image. Moreover, the coefficients of every matrix in the image is contained in the finitely generated integral domain $\mathbb{Q}[q^{\pm 1}, s^{\pm 1}]$ over $\mathbb{Q}$.

Choose an open subset $U \subseteq \mathrm{spec}(\mathbb{Q}[q^{\pm 1}, s^{\pm 1}])$ as in Lemma \ref{Open subset of good specializations}, such that at every closed point $x \in U$, the image of $F_2$ in $\mathrm{PGL}_N(k_x)$ is either finite or Zariski-dense. Moreover, by Corollary \ref{Extension of representation to AutF2}, we get an automorphism $\varphi \in \mathrm{Aut}(\mathrm{PSL}_N(\mathbb{K}))$ such that $\varphi\sigma_i\varphi^{-1} = \sigma_i^{-1}$, $\varphi^2 = 1$. Let $V \subseteq \mathrm{spec}(\mathbb{Q}[q^{\pm 1}, s^{\pm 1}])$ be an open subset where this automorphism is defined, obtained by inverting all denominators appearing in $\varphi, \varphi^{-1}$.

Pick a point $(q_0, s_0) \in \mathbb{Q}^2 \cap U \cap V,$ and further assume the open algebraic condition $q_0^as_0^b \neq \pm 1$. The image of $F_2$ at the residue field at $(q_0, s_0)$ is not finite - as $q_0^as_0^b \neq \pm 1$ is rational, and therefore not a root of unity. Hence, the image of $F_2$ is Zariski-dense.

After this specialization, we get a map $B_4 \to \mathrm{GL}_N(\mathbb{Q})$, such that the image of $F_2$ is contained in $\mathrm{SL}_N(\mathbb{Q})$, and $F_2$ has Zariski dense image in $\mathrm{PGL}_N(\mathbb{Q})$. Denote the image of $F_2$ by $\Gamma \leq \mathrm{SL}_N(\mathbb{Q}),$ and let $H$ denote its Zariski-closure in $\mathrm{SL}_N.$ The image of $H$ under the projection $\pi:\mathrm{SL}_N \to \mathrm{PGL}_N$ is all of $\mathrm{PGL}_N$. Since $\pi$ has finite fibers, $\dim(H)=\mathrm{dim}(\mathrm{PGL}_N)=\mathrm{dim}(\mathrm{SL}_N)$. However, as $\mathrm{SL}_N$ is simple, it has no proper closed subgroups of the same dimension, implying that $H = \mathrm{SL}_N$.

By the Strong Approximation Theorem, it follows that there are infinitely many specializations of the representation of $F_2$ to prime fields that map onto $\mathrm{SL}_N(\mathbb{F}_p)$ and in particular map onto $\mathrm{PSL}_N(\mathbb{F}_p).$

We have now shown that there are infinitely many specializations to finite fields such that $F_2 \to \mathrm{PSL}_N(\mathbb{F}_p)$ is surjective. This map extends to a map $B_4 \to \mathrm{PGL}_N(\mathbb{F}_p)$, and moreover as $(q_0, s_0) \in V$ the automorphism $\varphi$ restricts to an automorphism $\varphi \in \mathrm{Aut}(\mathrm{PSL}_N(\mathbb{Q}))$ such that $\varphi\sigma_i\varphi^{-1} = \sigma_i^{-1}$, $\varphi^2 = 1$. Further choosing the prime $p$ such that it does not appear in any of the denominators of $\varphi, \varphi^{-1}$, we may assume that $\varphi$ restricts to an automorphism in $\mathrm{Aut}(\mathrm{PSL}_N(\mathbb{F}_p)).$ By Lemma \ref{semi-direct product AutF2} this gives a map $\mathrm{Aut}(F_2) \to \mathrm{Aut}(\mathrm{PSL}_N(\mathbb{F}_p)).$ It now follows that the kernel of the map $F_2 \to \mathrm{PSL}_N(\mathbb{F}_p)$ is characteristic.
\end{proof}

\end{document}
